% =============================================================
% 03_development.tex — システム開発とプロジェクトマネジメント
% =============================================================
\chapter{システム開発とプロジェクトマネジメント}

\chaptermeta{%
  \item 開発工程の5段階を順序どおり説明できる
  \item ウォーターフォールとアジャイルの違いを即答できる
  \item QCD（品質・コスト・納期）のトレードオフを理解できる
  \item WBS/ガントチャートを使い分けられる
  \item テスト工程のV字モデルを説明できる
}{75分}{2}{第1--2章}

マネジメント系の前半は「つくる側」の話です。
どんな順番で開発し、どう管理するか——
\textbf{開発モデルの識別}と\textbf{PMの基本用語}を押さえれば得点源になります。

% =============================================================
\section{開発工程の5段階}
% =============================================================

\begin{teigi}[共通フレーム（SLCP）の5工程]
\begin{enumerate}
  \item \termtag{要件定義}：「何を作るか」を決める（ユーザ主体）
  \item \termtag{外部設計（基本設計）}：画面・帳票・インターフェースを決める
  \item \termtag{内部設計（詳細設計）}：プログラムの内部構造を決める
  \item \termtag{プログラミング}：実際にコードを書く
  \item \termtag{テスト}：単体→結合→システム→運用テストの順に実施
\end{enumerate}
\end{teigi}

\begin{pitfall}[要件定義と外部設計を混同しがち]
\begin{itemize}
  \item \textbf{要件定義}は「何をしたいか」を洗い出す（What）
  \item \textbf{外部設計}は「どう見せるか」を決める（How の外側）
  \item 要件定義はユーザ主体、外部設計は開発者主体
\end{itemize}
\end{pitfall}

\subsection{テスト工程とV字モデル}

\begin{hikaku}[V字モデル：設計とテストの対応]
\comptable{設計工程 & 対応するテスト & 確認内容}{
要件定義 & 運用テスト（承認テスト） & ユーザの要件を満たすか \\
外部設計 & システムテスト & システム全体が正しく動くか \\
内部設計 & 結合テスト & モジュール間の連携 \\
プログラミング & 単体テスト & 個々のモジュールの正しさ
}
\end{hikaku}

% =============================================================
\section{開発モデルの比較}
% =============================================================

\begin{hikaku}[代表的な開発モデル]
\comptable{モデル & 特徴 & 向いている場面}{
ウォーターフォール & 工程を順に進める・後戻りしにくい & 要件が明確な大規模開発 \\
アジャイル & 短い反復で少しずつ完成 & 要件変更が多いプロジェクト \\
プロトタイピング & 試作品で要件を確認 & ユーザのイメージが曖昧 \\
スパイラル & 設計→試作→評価を繰り返す & リスクの高い大規模開発 \\
DevOps & 開発と運用の一体化 & 継続的デリバリーが必要
}
\end{hikaku}

\begin{teigi}[アジャイルのキーワード]
\begin{itemize}
  \item \termtag{イテレーション}：1〜4週間の短い開発サイクル
  \item \termtag{スクラム}：代表的なアジャイルフレームワーク
  \item \termtag{スプリント}：スクラムにおけるイテレーション
  \item \termtag{プロダクトオーナー}：要件の優先順位を決める人
  \item \termtag{デイリースクラム}：毎日の短い進捗共有ミーティング
\end{itemize}
\end{teigi}

% --- 例題 ---
\begin{reidai}{開発モデルの識別}
システム開発において、短い期間で設計・実装・テストを繰り返し、
少しずつ機能を追加していく開発手法はどれか。

\sentakushi{ウォーターフォールモデル}{アジャイル開発}{共通フレーム}{PMBOK}
\end{reidai}

\begin{shishin}
「短い期間」「繰り返し」「少しずつ機能を追加」がキーワード。
ウォーターフォールは「順に進める」、共通フレームは「開発プロセスの規格」、
PMBOKは「PM知識体系」なので該当しない。
\end{shishin}

\begin{kaitou}
\textbf{正解：イ（アジャイル開発）}

短いイテレーションで反復的に開発する手法がアジャイル開発。
ウォーターフォールは工程を順に進め、後戻りしにくい点が対照的。
\end{kaitou}

\begin{minuki}[見抜き方]
\begin{itemize}
  \item 「繰り返し」「反復」「短期間」→ アジャイル
  \item 「順番に」「後戻りしない」「前工程完了後に次へ」→ ウォーターフォール
  \item 「試作品を見せて確認」→ プロトタイピング
  \item 「開発と運用の一体化」「CI/CD」→ DevOps
\end{itemize}
\end{minuki}

% =============================================================
\section{プロジェクトマネジメントとQCD}
% =============================================================

\begin{teigi}[QCD：プロジェクトの三大制約]
\begin{itemize}
  \item \termtag{Q（Quality）}：品質——成果物の正しさ・完成度
  \item \termtag{C（Cost）}：コスト——予算・工数
  \item \termtag{D（Delivery）}：納期——スケジュール
\end{itemize}
3つはトレードオフの関係にあり、1つを優先すると他が犠牲になりやすい。
\end{teigi}

\begin{pitfall}[「ファンクションポイント法」と「COCOMO」の混同]
\begin{itemize}
  \item \textbf{ファンクションポイント法}：機能の数と複雑さで工数を見積もる
  \item \textbf{COCOMO}：プログラムの行数で工数を見積もる
  \item 「機能」と「行数」のどちらに着目するかが違う
\end{itemize}
\end{pitfall}

% =============================================================
\section{WBS・ガントチャート}
% =============================================================

\begin{teigi}[WBSとガントチャート]
\begin{itemize}
  \item \termtag{WBS}（Work Breakdown Structure）：作業を階層的に分解した構造図。プロジェクトの全作業を洗い出す。
  \item \termtag{ガントチャート}：横軸に時間、縦軸に作業項目をとった棒グラフ。進捗管理に使う。
  \item \termtag{アローダイアグラム}：作業の依存関係を矢印で表す。\textbf{クリティカルパス}（最長経路）を見つけて全体の最短所要日数を求める。
\end{itemize}
\end{teigi}

\begin{hikaku}[PM手法の使い分け]
\comptable{ツール & 目的 & 特徴}{
WBS & 作業の洗い出し & ツリー構造で漏れなく分解 \\
ガントチャート & 進捗の見える化 & 横棒で期間を表示 \\
アローダイアグラム & 依存関係と最短日数 & クリティカルパスの特定 \\
PERT & 所要時間の見積もり & 楽観/最頻/悲観の3点見積もり
}
\end{hikaku}

% =============================================================
\section{過去問ドリル}
% =============================================================

\shoumatsuhead{過去問ドリル：開発とPM}

\begin{itpmon}[\sourcebadge{R6}\enspace\diffbadge{標}]{%
業務プロセスを可視化する手法のうち、業務の流れを「イベント」「アクティビティ」「ゲートウェイ」などの記号で表記するものはどれか。}
\sentakushi
  {BPMN（Business Process Model and Notation）}
  {DFD（Data Flow Diagram）}
  {E-R図}
  {UML}
\end{itpmon}
\begin{kaisetsu}{ア}
\textbf{BPMN}は業務プロセスをイベント・アクティビティ・ゲートウェイの記号で可視化する標準表記法。DFDはデータの流れ、E-R図はデータの関係、UMLはオブジェクト指向の設計図。
\end{kaisetsu}

\begin{itpmon}[\sourcebadge{original}\enspace\diffbadge{標}]{%
アジャイル開発の特徴として、最も適切なものはどれか。}
\sentakushi
  {全工程の計画を詳細に立ててから開発に着手する。}
  {短い期間での反復を繰り返し、段階的にシステムを完成させる。}
  {利用者にプロトタイプを見せて要件を確定する。}
  {プログラムのステップ数から開発コストを見積もる。}
\end{itpmon}
\begin{kaisetsu}{イ}
アジャイルは短いイテレーションで反復的に開発する手法。アはウォーターフォール、ウはプロトタイピング、エはCOCOMOの説明。
\end{kaisetsu}

\begin{itpmon}[\sourcebadge{original}\enspace\diffbadge{標}]{%
WBS（Work Breakdown Structure）の説明として、最も適切なものはどれか。}
\sentakushi
  {プロジェクトの作業を階層的に分解して構造化したもの}
  {横軸に時間、縦軸に作業をとって進捗を管理する図表}
  {作業の依存関係を矢印で表し、クリティカルパスを求める手法}
  {品質・コスト・納期の三大制約を管理するフレームワーク}
\end{itpmon}
\begin{kaisetsu}{ア}
\textbf{WBS}はプロジェクトの全作業を階層的に分解した構造。イはガントチャート、ウはアローダイアグラム（PERT図）、エはQCDの説明。
\end{kaisetsu}

% =============================================================
\section{タイムアタック}
% =============================================================

\begin{timeattack}{開発・PMクイックチェック}{5分}{4問中3問}

\begin{itpmon}[\diffbadge{基}]{システム開発の工程で、最初に行うのはどれか。}
\sentakushi{外部設計}{テスト}{プログラミング}{要件定義}
\end{itpmon}
\begin{kaisetsu}{エ}
開発は\textbf{要件定義}→外部設計→内部設計→プログラミング→テストの順。
\end{kaisetsu}

\begin{itpmon}[\diffbadge{基}]{V字モデルで、要件定義に対応するテストはどれか。}
\sentakushi{単体テスト}{結合テスト}{システムテスト}{運用テスト}
\end{itpmon}
\begin{kaisetsu}{エ}
要件定義 ↔ \textbf{運用テスト}。外部設計 ↔ システムテスト、内部設計 ↔ 結合テスト、プログラミング ↔ 単体テスト。
\end{kaisetsu}

\begin{itpmon}[\diffbadge{標}]{スクラムにおける「スプリント」の説明として適切なものはどれか。}
\sentakushi
  {プロジェクト全体の工程を順に進める方式}
  {1〜4週間の短い反復開発の単位}
  {開発と運用を一体化する取組み}
  {試作品を使って要件を確認する手法}
\end{itpmon}
\begin{kaisetsu}{イ}
\textbf{スプリント}はスクラムにおける1〜4週間のイテレーション。アはウォーターフォール、ウはDevOps、エはプロトタイピング。
\end{kaisetsu}

\begin{itpmon}[\diffbadge{標}]{プロジェクトの全体所要日数を求めるために使う「クリティカルパス」とは何か。}
\sentakushi
  {コストが最も大きい作業の経路}
  {品質リスクが最も高い作業の経路}
  {所要時間が最も長い作業の経路}
  {担当者が最も多い作業の経路}
\end{itpmon}
\begin{kaisetsu}{ウ}
\textbf{クリティカルパス}はアローダイアグラムにおける最長経路であり、プロジェクト全体の最短所要日数を決定する。
\end{kaisetsu}

\end{timeattack}

% =============================================================
\section{よくあるミスと到達確認}
% =============================================================

\begin{yokumiss}[この章でよくあるミス]
\begin{enumerate}[leftmargin=1.6em, itemsep=5pt]
  \item \textbf{V字モデルの対応関係を間違える}——「一番上同士、一番下同士」で対応する。
  \item \textbf{アジャイルを「計画しない開発」と誤解する}——計画はする。ただし短いサイクルで柔軟に変更する。
  \item \textbf{WBSとガントチャートを混同する}——WBSは「分解」、ガントチャートは「日程管理」。
  \item \textbf{ファンクションポイント法とCOCOMOの違いを忘れる}——「機能数」か「行数」か。
  \item \textbf{クリティカルパスを「最短経路」と勘違いする}——最長経路が正解。最長経路＝全体の最短所要日数。
\end{enumerate}
\end{yokumiss}

\begin{tatsaku}
  \item 開発工程の5段階を順番に言える
  \item V字モデルの対応関係を4組すべて言える
  \item ウォーターフォールとアジャイルの違いを3点以上挙げられる
  \item アジャイル関連用語（スクラム/スプリント/PO）を説明できる
  \item QCDの三大制約とトレードオフを説明できる
  \item WBS/ガントチャート/アローダイアグラムの目的を区別できる
  \item クリティカルパスの意味と求め方を説明できる
  \item BPMNとDFDの違いを説明できる
\end{tatsaku}

\nextchapter{%
サービスマネジメントと監査に進みます。SLA/SLM、ITIL、稼働率の計算、
システム監査の独立性——「運用する側」の視点を学びます。
}
